% titel: Blatt 7, A3: MO-Energiediagramm und Bindungsordnung von NO
% skript: Mortimer Kap. 10 (Molekülorbitale)

% Teil I: Werkzeuge, kein L\"osungsweg.

F\"ur ein zweiatomiges Molek\"ul AB tr\"agt man die
Molek\"ulorbital-Energieniveaus (aus den 2s- und 2p-AO der Atome) auf und
f\"ullt sie mit den Valenzelektronen auf.

\begin{werkzeugkasten}[MO-Aufbau und Bindungsordnung]
\begin{itemize}
  \item Reihenfolge (f\"ur $Z\geq 8$, z.~B. NO wie O$_2$):
        \begin{center}
        $\sigma_{2s} < \sigma^*_{2s} < \pi_{2p} \;(\pi_x,\pi_y)<
         \sigma_{2p} < \pi^*_{2p} \;(\pi^x_*,\pi^y_*)$
        \end{center}
  \item 2 Valenzelektronen pro Niveau (Pauli), erst einzeln (Hund).
  \item Bindungsordnung: $(2\times\text{$\pi$\,(besetzte bindende)}+
        \text{besetzte $\sigma$}) - (2\times\text{antibindende})$, geteilt
        durch 2.
  \item Ungepaarte Elektronen $\rightarrow$ paramagnetisch.
\end{itemize}
\textit{Skript:} Mortimer Kap. 10, Molek\"ulorbitale bei zweiatomigen
Molek\"ulen.
\end{werkzeugkasten}

\begin{achtung}[Darauf achten]
\begin{itemize}
  \item NO hat eine ungerade Zahl an Valenzelektronen ($5+6 = 11$)
        $\rightarrow$ ein Elektron wird in ein
        $\pi^*_{2p}$-Niveau gelegt: ungepaart.
  \item Das antibindende $\pi^*$-Elektron senkt die Bindungsordnung.
  \item NO$^+$ hat ein Elektron weniger (10) als NO (11).
\end{itemize}
\end{achtung}
